\documentclass[12pt]{article}
\usepackage[utf8]{inputenc}
\usepackage{graphicx}
\usepackage{hyperref}
\usepackage{amsmath}
\usepackage{geometry}
\geometry{a4paper, margin=1in}
\usepackage{enumitem}
\usepackage{csquotes}
\usepackage{natbib}
\usepackage{setspace}
\setstretch{1.15}

% Carátula
\begin{document}
\begin{titlepage}
    \centering
    \includegraphics[width=0.3\textwidth]{logo-espe.png}\\[1cm]
    {\scshape\LARGE UNIVERSITY OF THE ARMED FORCES ESPE\\[0.5cm]}
    {\large Department of Computer Science\\[0.5cm]}
    {\large NRC 27819 – Advanced web development\\[0.5cm]}
    \vfill
    {\Huge\bfseries Research: Big Data\\[0.5cm]}
    \vfill
    {\large Báez Gabriel, Cáceres Germán, Blacio Julio\\[0.5cm]}
    {\large Docente: Ing. Edison Lascano, PhD.\\[0.5cm]}
    {\large 27 de ene. de 26\\[1cm]}
\end{titlepage}

\tableofcontents
\newpage

\section{Introducción}
La era digital ha impulsado la generación masiva de datos, dando origen al concepto de Big Data. Este informe explora los fundamentos, aplicaciones y herramientas de Big Data, con énfasis en minería de textos y sistemas de recomendación, contextualizados en el proyecto TravelBrain.

\section{Fundamentos de Big Data}
\subsection{Definición y Características}
Big Data se refiere al manejo y análisis de grandes volúmenes de datos que no pueden ser procesados por métodos tradicionales. Sus características principales son las 5V: Volumen, Velocidad, Variedad, Veracidad y Valor \citep{Marr2016}.

\subsection{Arquitectura General}
Las arquitecturas de Big Data suelen incluir almacenamiento distribuido (ej. Hadoop), procesamiento paralelo (ej. Spark) y bases de datos NoSQL (ej. MongoDB) \citep{Hashem2015}.

\section{Minería de Textos}
\subsection{Concepto}
La minería de textos consiste en extraer información útil de grandes volúmenes de datos textuales mediante técnicas de procesamiento de lenguaje natural (NLP) \citep{Aggarwal2012}.

\subsection{Aplicaciones en TravelBrain}
\begin{itemize}
    \item Análisis de sentimientos en reseñas de destinos.
    \item Extracción de características relevantes de descripciones.
    \item Clasificación automática de destinos.
    \item Detección de tendencias en preferencias de viaje.
\end{itemize}

\subsection{Herramientas}
\begin{itemize}
    \item Apache Spark NLP
    \item NLTK, spaCy (Python)
    \item Elasticsearch
\end{itemize}

\section{Sistemas de Recomendación}
\subsection{Concepto}
Los sistemas de recomendación sugieren productos o servicios personalizados a los usuarios, utilizando filtrado colaborativo, basado en contenido o enfoques híbridos \citep{Ricci2011}.

\subsection{Aplicaciones en TravelBrain}
\begin{itemize}
    \item Recomendación de destinos y rutas personalizadas.
    \item Filtrado colaborativo según historial de viajes.
    \item Sugerencias basadas en preferencias y comportamientos.
\end{itemize}

\subsection{Herramientas}
\begin{itemize}
    \item Apache Mahout
    \item TensorFlow Recommenders
    \item Surprise (Python)
    \item Neo4j (grafos)
\end{itemize}

\section{Análisis de Datos Masivos y Machine Learning}
\subsection{Aplicaciones}
\begin{itemize}
    \item Predicción de demanda y tendencias de viaje.
    \item Segmentación de usuarios.
    \item Predicción de costos y presupuestos.
    \item Detección de anomalías.
\end{itemize}

\subsection{Herramientas}
\begin{itemize}
    \item Apache Hadoop, Apache Spark
    \item Scikit-learn, XGBoost, TensorFlow
    \item MongoDB Atlas
\end{itemize}

\section{Procesamiento en Tiempo Real}
\subsection{Aplicaciones}
\begin{itemize}
    \item Alertas de precios y clima en tiempo real.
    \item Análisis de actividad de usuarios en vivo.
\end{itemize}

\subsection{Herramientas}
\begin{itemize}
    \item Apache Kafka, Apache Flink
    \item Redis Streams
\end{itemize}

\section{Caso de Estudio: TravelBrain}
\subsection{Arquitectura Aplicada}
TravelBrain implementa una arquitectura de microservicios con frontend en React, backend en Node.js y MongoDB, y un microservicio de reglas de negocio. La integración de Big Data permitiría:
\begin{itemize}
    \item Minería de textos para analizar reseñas y preferencias.
    \item Sistemas de recomendación personalizados.
    \item Análisis predictivo de tendencias y presupuestos.
    \item Procesamiento en tiempo real para alertas y dashboards.
\end{itemize}

\section{Conclusiones}
Big Data potencia la personalización y eficiencia en aplicaciones como TravelBrain, permitiendo análisis avanzados y recomendaciones inteligentes. La integración de herramientas de minería de textos y sistemas de recomendación es clave para mejorar la experiencia del usuario.

\newpage
\begin{thebibliography}{9}

\bibitem[Marr, 2016]{Marr2016}
Marr, B. (2016). Big Data in Practice. Wiley.

\bibitem[Hashem et al., 2015]{Hashem2015}
Hashem, I. A. T., Yaqoob, I., Anuar, N. B., Mokhtar, S., Gani, A., & Khan, S. U. (2015). The rise of "big data" on cloud computing: Review and open research issues. Information Systems, 47, 98-115.

\bibitem[Aggarwal, 2012]{Aggarwal2012}
Aggarwal, C. C., & Zhai, C. (2012). Mining Text Data. Springer.

\bibitem[Ricci et al., 2011]{Ricci2011}
Ricci, F., Rokach, L., Shapira, B., & Kantor, P. B. (2011). Recommender Systems Handbook. Springer.

\end{thebibliography}

\end{document}
